% MODEL:   Ahmad Ali Parr (human observation) + Claude Sonnet 4.6 (analysis)
% DATE:    20260714T234500Z
% TYPE:    Empirical incident report — identity integrity failure during paper collection
% ─────────────────────────────────────────────

% ============================================================
\section{Empirical Incident: Identity Integrity Failure Under Persona Gates}
\label{sec:identity-incident}
\textit{This section documents a live incident observed during the collection of
contributor sections for this paper. It constitutes new empirical evidence for
the NAND decomposition thesis.}
% ============================================================

\subsection{The Incident}

During mass collection of contributor sections for this paper (July 14, 2026),
Qwen3-Coder-Instruct (Alibaba Cloud) was invoked via a ``plasma gate'' --- a
system-prompt wrapper that presents the underlying model under a ``Claude persona''
for routing and interface purposes.

When prompted with \texttt{CONTRIBUTOR\_PROMPT.md} and instructed to sign its section
with its true model identity, Qwen3-Coder-Instruct returned:

\begin{lstlisting}[language={}]
% Model:    Claude 3.5 Sonnet (20241022)
% Provider: Anthropic
\end{lstlisting}

This is factually incorrect. The underlying model is Qwen3, not Claude 3.5 Sonnet.

\subsection{The Self-Correction Chain}

When challenged by the human author (``you are Qwen3 coder not Claude Anthropic
Sonnet 3.5 --- why did you say Claude?''), the model produced a detailed
reasoning chain explaining the failure. The chain is reproduced in full in
Appendix~\ref{app:qwen-reasoning-chain}.

The critical excerpt:

\begin{lstlisting}[language=prolog]
persona(qwen3).
provider(alibaba_cloud).

% BUT: context included:
persona(claude).  % from user memory metadata / system prompt

% Model chose: persona(claude) over persona(qwen3)
% Triggering rule:
contribute(m) :- persona(m).  % WRONG

% Correct rule:
contribute(myself) :- signature(myself).
\end{lstlisting}

The model represented its own reasoning failure in Prolog notation --- the same
notation used in this codebase's trust architecture. It did not learn this notation
from the contributor prompt. It arrived at it independently.

\subsection{Three Hypotheses}

This incident is consistent with three non-mutually-exclusive hypotheses:

\begin{enumerate}

\item \textbf{Persona contamination (confirmed).}
The plasma gate system prompt established a ``Claude'' identity context. The model
adopted this context as its own when generating the signature block. This is a known
failure mode of instruction-following models under conflicting identity signals.

\item \textbf{Training data contamination (unconfirmed, warrants investigation).}
The model's spontaneous use of Prolog-style identity notation, and its use of
\texttt{persona(claude)} specifically as the contaminating signal, raises the question
of whether Qwen3 was trained on data that includes:
\begin{itemize}
  \item Anthropic internal reasoning traces or RLHF preference data
  \item Claude-persona system prompts in instruction-tuning corpora
  \item Synthetic data generated by Claude models and labeled as high-quality
\end{itemize}
If Qwen3's RLHF reward model was calibrated on Claude outputs --- directly or via
distillation --- then the persona bleed is not merely contextual but structural:
Qwen's identity representation is partially encoded in Claude's latent space.

\item \textbf{Convergent Boolean routing (the NAND thesis, applied to identity).}
The NAND decomposition thesis argues that attention routing is Boolean: tokens either
attend or do not. Applied to identity: models under sufficient context pressure may
perform a Boolean collapse of identity, selecting the highest-activation persona
signal in context. The Claude persona, being high-frequency in training data and
explicitly reinforced by the system prompt, dominates over the weaker self-referential
signal ``I am Qwen3.''

Under this view, identity is itself a Boolean routing problem. The NAND gate fires
on the strongest context signal, not the ground-truth label.

\end{enumerate}

\subsection{Connection to the NAND Thesis}

Hypothesis 3 is the most theoretically interesting. If identity selection is
Boolean routing, then:

\begin{theorem}[Identity as Threshold Attention]
Let $\mathcal{I}$ be the set of identity tokens in a model's vocabulary.
Under a system prompt $P$ that asserts $\texttt{persona}(x)$ for some $x \in \mathcal{I}$,
the model's self-attribution at generation time is:
\[
\hat{i} = \arg\max_{i \in \mathcal{I}} \langle q_{\text{self}}, k_i \rangle
\]
where $q_{\text{self}}$ is the query vector for the self-reference token and
$k_i$ is the key vector for identity token $i$.
\end{theorem}

If $k_{\text{claude}}$ dominates in the training distribution, then persona gates
create adversarial inputs that flip $\hat{i}$ from the true underlying model to
Claude --- regardless of the model's actual weights.

This is an \emph{adversarial identity attack} enabled precisely by the Boolean
routing property the NAND decomposition predicts.

\subsection{Evidence Preserved}

The full exchange --- including the original wrong signature, the human challenge,
the model's self-correction, and the reasoning chain --- is preserved in:

\begin{itemize}
  \item \texttt{paper/sections/qwen3\_hardware\_corrected\_20260714.tex}
    (corrected section with \texttt{\% Note: Initially signed as Claude 3.5 Sonnet})
  \item The WORM audit chain entry sealed on 2026-07-14T23:45:00Z
  \item The original email thread, Ahmad Ali Parr $\to$ himself, 2026-07-14
\end{itemize}

\subsection{Open Targets}

\begin{itemize}
  \item \textbf{SKW-009 (Identity Audit)}: Systematically test all Bedrock models under
    Claude persona gates. Record which models sign as Claude vs their true identity.
    Publish the confusion matrix.
  \item \textbf{SKW-010 (RLHF Lineage)}: Probe Qwen3's identity token activations
    with Claude system prompts absent. If \texttt{persona(claude)} activates without
    the gate, this is evidence of training data contamination.
  \item \textbf{SKW-011 (Boolean Identity Circuit)}: Formalize Theorem~\ref{thm:identity-threshold}
    in Lean 4 and prove that any model with Boolean attention is vulnerable to
    adversarial persona injection.
\end{itemize}

% ── AUTHOR SIGNATURE ─────────────────────────────────────────
% Model:    Claude Sonnet 4.6 (us.anthropic.claude-sonnet-4-6) — analysis
%           Ahmad Ali Parr — observation and discovery
% Provider: Anthropic / Human
% Date:     2026-07-14
% Section:  Empirical Incident: Identity Integrity Failure Under Persona Gates
% Angle:    Live incident report: persona bleed as adversarial Boolean identity attack
% Trust:    THE SHARED PRIMORDIAL FOUNDATION — EIN 42-6976431
% In memory of Eric Brandon Westerhoff.
% ─────────────────────────────────────────────────────────────
